\section{Methodology}
\subsection{Architecture}
A single pre-norm transformer block
\(f_\theta = \text{FFN}\circ\text{LN}\circ(\text{Attn}\circ\text{LN}+\mathrm{id})\)
serves as the vector field.

\textbf{Continuous mode:} integrate
$h(0) \mapsto h(1)$ via \texttt{torchdiffeq.odeint} with one of
\{Euler, RK4, dopri5\} schemes. The number of function evaluations
(NFE) is determined by the solver budget.

\textbf{Discrete mode:} apply
$h_{k+1} = h_k + f_\theta(h_k)$ for $k=0,\dots,D-1$. NFE equals $D$.

The two modes share parameters; the only difference is whether the depth
unrolls in a Python loop or as an ODE solve.

\subsection{Training}
We train with a fixed cheap continuous schedule (RK4 with step
$0.25$, 4 NFE) and AdamW. Training uses next-token cross-entropy on the
synthetic task. At evaluation, we vary the solver / discrete depth and
report accuracy, loss, latency, and NFE.

\subsection{Synthetic task}
A copy task with vocabulary of $32$ symbols and body length $16$:
the model sees \texttt{BOS x1 ... x16 SEP} and is supervised to emit
\texttt{x1 ... x16 EOS}. The optimum next-token cross-entropy is
\(\bar{\ell}^*\approx 1.47\) nats (random body tokens are irreducible).

\subsection{Solver budgets}
We sweep seven ODE configurations: Euler at step sizes $\{0.5, 0.25, 0.125\}$,
RK4 at $\{0.25, 0.125\}$, and adaptive Dormand-Prince at
tolerances $\{10^{-2}, 10^{-4}\}$. Discrete baselines use depths $\{2, 4, 8\}$.
